\begin{frame}{할인율 $\gamma = 1$ 을 쓰는 이유}
  \begin{itemize}
    \item Week 4 에서는 $\gamma < 1$ 을 써서 ``무한히 길어질 수
      있는 미래 보상''을 수렴되게 했었다.
    \item 블랙잭의 한 판은 \textbf{결정적으로 끝나는 유한
      호라이즌} --- 누군가 버스트하거나 딜러가 17 이상으로
      마무리되는 순간 에피소드가 끝나므로, 보상이 ``무한히 먼
      미래''까지 이어지지 않는다.
    \item 그런 \textbf{유한하고 끝이 나는 과업}에서는
      $\gamma = 1$ (보상을 통째로 가중하지 않는 것)이
      자연스러운 선택.
  \end{itemize}
  \vskip0.5em
  \begin{block}{보상을 $-1, 0, +1$ 로 두는 이유}
    ``기대 수익이 0 근처의 게임''이라는 직관이 $Q$ 값의
    \textbf{부호와 크기}에 그대로 담기게 한다. $Q(s,a) > 0$
    은 기대적으로 \textbf{이길} 확률이 더 크다, $Q(s,a) < 0$
    은 \textbf{지실} 확률이 더 크다.
  \end{block}
\end{frame}
